%% ELEC2013 2007

%% ----------------------------------------------------------------
%% exam.tex
%% ----------------------------------------------------------------
\documentclass{sotonExamBoxes}    % use option [answers] to include them
%\documentclass[answers]{uosexamb}
%% ----------------------------------------------------------------
\usepackage{graphicx}
%\usepackage{pstricks}
\usepackage{alltt}
\usepackage{bm}
\usepackage{amssymb}
\usepackage{amsfonts}
\usepackage{amsmath}
\usepackage{stmaryrd}
\thinmuskip=1mu
%\renewcommand{\ttdefault}{pcr}
%% macros.tex -- shared macro definitions for the problem sheets.
%%
%% Every sheet inputs this instead of carrying its own copy. Input it after
%% the \usepackage lines: the definitions below use bm, amsmath and amssymb.
%%
%% Six macros were defined inconsistently across the sheets when this file
%% was made; the choices are marked CONFLICT below.




\usepackage{stmaryrd}

\renewcommand{\ttdefault}{pcr}

\newcommand{\tr}{\textsf{T}}

\newcommand{\e}[1]{{\rm e}^{#1}}

\newcommand{\E}[1]{{\rm e}^{{\displaystyle #1}}}

\newcommand{\logg}[1]{\log\left(#1\right)}

\newcommand{\Prob}[1]{\mathbb{P}\left(#1\right)}

\newcommand{\Step}{\mathop{\mathrm{O}\hspace{-15pt}\rule[8.5pt]{11pt}{1.2pt}\hspace{2pt}}}

\newcommand\diag{\mathop{\mathrm{diag}}}

\newcommand{\grad}{\bm{\nabla}}

\newcommand{\dd}{\mathrm{d}}

\newcommand{\class}{\mathcal{C}}

\newcommand{\data}{\mathcal{D}}

%% CONFLICT: gaussianProcess.tex had {U}{eur}{b}{n} (Euler bold);
%% every other sheet had this. Matrices in that one sheet now look
%% different -- see the report.

\DeclareMathAlphabet{\mat}{OT1}{cmss}{bx}{n}

%% The identity needs to be distinguishable. In cmss a bold I is a plain
%% vertical bar with no crosspieces, so \mat{I} reads as a rule rather than a
%% letter. Euler's bold I keeps its serifs, so the identity alone is set in
%% Euler while everything else stays cmss.
\DeclareMathAlphabet{\matEuler}{U}{eur}{b}{n}
\newcommand{\Id}{\matEuler{I}}

%% CONFLICT and BUG FIX: 16 sheets defined this using \qpartial,
%% which is defined nowhere, so \pd was broken in all of them.
%% Only problem2.tex had the working \partial form, kept here.

\newcommand{\pd}[2]{\ensuremath{\frac{\partial #1}{\partial #2}}\xspace}

\newcommand{\M}[1]{\ensuremath{\boldsymbol{#1}}\xspace}

\newcommand{\bst}{\rule{0mm}{5mm}}

%% CONFLICT: some sheets omitted the \! . Kept the tighter spacing,
%% which is what every sheet that actually uses \av had.

\newcommand{\av}[2][\,]{\mathbb{E}_{#1}\!\left[ {\strut #2} \right]}

\newcommand{\len}[1]{\| #1 \|}

\newcommand{\inner}[2]{\left\langle #1, #2\right\rangle}

\newcommand{\squeeze}{\setlength{\itemsep}{-2pt}}

%% CONFLICT: two spacings, identical output in maths mode.

\newcommand{\bra}[1]{\left( #1 \right)}

\newcommand{\pred}[1]{\left\llbracket { \small #1} \right\rrbracket}

\newcommand{\mln}{\hat{f}_n(\bm{x})}


\newcommand{\Tr}{\mathop{\mathrm{tr}}\,}

%% NOT SHARED: gaussianProcess.tex redefines \det to take an argument
%% (\det{\mat{V}}), while calculus.tex and vectorSpaces.tex use the standard
%% no-argument operator (\det(\mat{X})). The two cannot coexist, so that
%% redefinition stays local to gaussianProcess.tex.

%% NOT SHARED: \ml and \mm mean different things in different sheets.
%% amgm and convexity use a no-argument \ml and \mm = \hat{m}_c(x);
%% ensembleLearning uses a one-argument \ml and \mm = \hat{m}(x).
%% Both forms are in use, so each sheet keeps its own.

%% local: this sheet's own meaning for \ml and \mm
\newcommand{\ml}{\hat{f}_c(\bm{x}|\bm{\theta})}
\newcommand{\mm}{\hat{m}_c(\bm{x})}
\graphicspath{{figures/} {../2018-19/figures}}
\graphicspath{{figures/}}


% \usepackage{picture}
\begin{document}
\unitTitle{Convexity and SVMs Problem Sheet}
\authors{Adam Pr\"ugel-Bennett}
\unitcode{AICE3001/AICE6002}
\semester{2}
\year{2022}
\duration{45}{0}


\maketitle

This sheet covers convexity, support vector machines and kernels.
Questions~1 and~2 are about convex functions and Jensen's inequality.
Question~3 solves a tiny SVM by hand, in both the primal and the dual
form.  Question~4 asks which functions are valid kernels.

% ***************************** question 1 **************************


\begin{question}{10}
  \begin{qparts}
    \qpart[4]{Starting from the definition of a convex function
      \begin{align}\label{eq:convex1}
        f\bra{a\,x + (1-a)\,y} \leq a\,f(x) + (1-a) \, f(y)
      \end{align}
      Let $a = \epsilon/(x-y)$ and rearrange the
      inequality to give
      \begin{align*}
        (x-y)\,\bra{\frac{f(y+\epsilon) - f(y)}{\epsilon}}
      \end{align*}
      on the left-hand side.  Taking the limit $\epsilon \rightarrow 0$
      show that the function $f(x)$ lies above the tangent line
      $t(x) = f(y) + (x-y)\,f'(y)$ going through the point $y$.}{\lines{12}}
      \begin{answer}
        Rearranging the Equation~\ref{eq:convex1}
        \begin{align*}
          f\bra{y + a(x-y)} \leq f(y) + a\,\bra{f(x) - f(y)}.
        \end{align*}
        Or
        \begin{align*}
          \frac{f\bra{y + a(x-y)} - f(y)}{a} \leq f(x) - f(y).
        \end{align*}
        Letting $a =\epsilon/(x-y)$ then
        \begin{align*}
          (x-y) \,  \frac{f\bra{y + \epsilon} - f(y)}{\epsilon} \leq
          f(x) - f(y)
        \end{align*}
        Taking the limit $\epsilon \rightarrow 0$ then using
        \begin{align*}
          \lim_{\epsilon \rightarrow 0} \frac{f\bra{y + \epsilon} -
          f(y)}{\epsilon}  = f'(y)
        \end{align*}
        So that
        \begin{align*}
           (x-y) \, f'(y) \leq f(x) - f(y)
        \end{align*}
        or
        \begin{align*}
          f(x) \geq f(y) + (x-y)\,f'(y) = t(x).
        \end{align*}
      \end{answer}
      \clearpage
      \qpart[1]{Sketch the tangent line, $t(x)$, at the point $y$ in the graph
        shown below.}{\qfig[7cm]{tangentline_q}}
      \begin{answer}
        \afig[7cm]{tangentline_a}
      \end{answer}
      \qpart[4]{Starting from the inequality for a convex function
        \begin{align}\label{eq:convex2}
          f(x) \geq f(y)  + (x-y)\, f'(y)
        \end{align}
        consider the case $y = x+\epsilon$, then by Taylor expanding $f(x+\epsilon)$
        and $f'(x+\epsilon)$ around $x$ and keeping all terms up to
        order $\epsilon^2$ show that for a convex function $f''(x)\geq0$.
      }{\lines{8}}
      \begin{answer}
        We use the expansions
        \begin{align*}
          f(x+\epsilon) &= f(x) + \epsilon\,f'(x) +
                          \frac{\epsilon^2}{2} f''(x) + O(\epsilon^3) \\
          f'(x+\epsilon) &= f'(x) + \epsilon\,f''(x) +
                           O(\epsilon^2) 
        \end{align*}
        Using $y=x+\epsilon$ and substituting into the Equation~(\ref{eq:convex2})
        \begin{align*}
          f(x) &\geq f(x+\epsilon)  - \epsilon \, f'(x+\epsilon) \\
          f(x) &\geq f(x) + \epsilon\,f'(x) +
          \frac{\epsilon^2}{2} f''(x) + O(\epsilon^3)
          - \epsilon\,\bra{f'(x) + \epsilon\,f''(x) +
                           O(\epsilon^2) }
        \end{align*}
        Or subtracting $f(x)$ from both sides
        \begin{align*}
          0 \geq - \frac{\epsilon^2}{2} \, f''(x) + O(\epsilon^3)
        \end{align*}
        Since this has to be true for all $\epsilon>0$ this requires $f''(x)\geq0$.
      \end{answer}
      \qpart[1]{Prove that $-\log(x)$ is convex-up for $x>0$.}{\lines{1}}
      \begin{answer}
        \begin{align*}
          \frac{\dd^2 (-\log(x))}{\dd \, x^2} = \frac{1}{x^2} \geq 0
        \end{align*}
      \end{answer}
  \end{qparts}
\end{question}
\clearpage


% ***************************** question 2 **************************


\begin{question}{10}
  \begin{qparts}
    \qpart[5]{If $\len{\bm{x}}$ is a proper norm use the triangular
      inequality
      ($\len{\bm{x} + \bm{y}}\leq \len{\bm{x}} + \len{\bm{y}}$),
      ``linearity'' of a norm ($\len{a\,\bm{x}} = |a|\,\len{\bm{x}}$) and
      the definition of convexity, to show that the norm is
      convex.}{\lines{12}}
    \begin{answer}
      For any vectors, $\bm{x}$, and $\bm{y}$ and any scalar
      $a\in[0,1]$ then
      \begin{align*}
        \len{a\, \bm{x} + (1-a)\, \bm{y}} \leq \len{a\,\bm{x}} +
        \len{(1-a)\,\bm{y}} = a\, \len{\bm{x}} + (1-a)\, \len{\bm{y}}
      \end{align*}
      where the first inequality follows from the triangular
      inequality and the second equality from the ``linearity'' of the
      norm (with $a,1-a\geq0$).  However,
      \begin{align*}
         \len{a\, \bm{x} + (1-a)\, \bm{y}} \leq a\, \len{\bm{x}} + (1-a)\, \len{\bm{y}}
      \end{align*}
      is the defining equation of convexity.
    \end{answer}
    \clearpage


     \qpart[5]{Consider a classification problem where $\ml$ is the
       probability that a learning machine with parameters
       $\bm{\theta}$ predicts that input $\bm{x}$ belongs to class
       $c\in\mathcal{C}$.  Assume the training is stochastic so the
       probability of obtaining parameters $\bm{\theta}$ is
       $\rho(\bm{\theta})$. Let $\mm = \av[\bm{\theta}]{\ml}$ be the output
       of the mean machine for class $c$.  Assuming that
       for a data point $(\bm{x}, y)$, where $y$ is a
       class label, we use a cross entropy loss
       \begin{align*}
         L(\bm{x}, y, \bm{\theta}) = -\sum_{c \in \mathcal{C}}
         \pred{y=c} \logg{\ml},
       \end{align*}
       show that  the expected loss over inputs and parameters can be
       written as the expected loss of the mean machine plus a second
       loss.  Use Jensen's inequality ($\av{\log(X)}\leq\logg{\av{X}}$) to
       show the second term is non-negative.}{\lines{15}}
     \begin{answer}
       \begin{align*}
         \bar{L} &= \av[(\bm{x},y)]{\av[\bm{\theta}]{  -\sum_{c \in \mathcal{C}}
                       \pred{y=c} \logg{\ml}}}
         \\
         &=  -\av[(\bm{x},y)]{\sum_{c \in \mathcal{C}} \pred{y=c} \logg{\mm}}
           - \av[(\bm{x},y)]{\av[\bm{\theta}]{  \sum_{c \in \mathcal{C}}
                       \pred{y=c} \logg{\frac{\ml}{\mm}}}}
       \end{align*}
       The first term acts like a bias.  The second  (variance-like) term is
       \begin{align*}
         - \av[(\bm{x},y)]{ \sum_{c \in \mathcal{C}}
         \pred{y=c} \av[\bm{\theta}]{ \logg{\ml}}}
         + \av[(\bm{x},y)]{\sum_{c \in \mathcal{C}} \pred{y=c} \logg{\mm}}
       \end{align*}
       But using Jensen's inequality $\av[\bm{\theta}]{ \logg{\ml}}
     \leq \logg{\av[\bm{\theta}]{\ml}} = \logg{\mm}$. Thus this second
     term is non-negative.
     \end{answer}
  \end{qparts}
\end{question}
\clearpage

% ***************************** question 3 **************************


\begin{question}{10}
  In the lectures the maximum-margin SVM finds $\bm{w}$ and $b$ that
  minimise $\tfrac12\len{\bm{w}}^2$ subject to
  $y_k\,(\bm{w}^\tr\bm{x}_k - b)\geq 1$ for every data point.  The dual
  problem is to maximise
  \begin{align*}
    \mathcal{L}(\bm{\alpha}) = \sum_k \alpha_k - \frac12 \sum_{k,l}
    \alpha_k\,\alpha_l\,y_k\,y_l\,\bm{x}_k^\tr\bm{x}_l
  \end{align*}
  subject to $\sum_k \alpha_k\,y_k = 0$ and $\alpha_k\geq0$, after which
  $\bm{w} = \sum_k \alpha_k\,y_k\,\bm{x}_k$.  In this question we solve
  a tiny problem by hand.
  \begin{qparts}
    \qpart[2]{Take two points, $\bm{x}_1 = (1,1)^\tr$ with $y_1=+1$ and
      $\bm{x}_2 = (-1,-1)^\tr$ with $y_2 = -1$.  By symmetry the
      solution has $b=0$ and $\bm{w} = (u,u)^\tr$.  Find $u$ and the
      margin $1/\len{\bm{w}}$.  Check that the margin is half the
      distance between the two points.
      \begin{answer}
        Both constraints read $2u\geq1$, and we want $u$ as small as
        possible, so $u=\tfrac12$ and $\bm{w} = (\tfrac12,\tfrac12)^\tr$.
        The margin is $1/\len{\bm{w}} = 1/(1/\sqrt2) = \sqrt2$.  The
        points are $\len{\bm{x}_1-\bm{x}_2} = \len{(2,2)^\tr} = 2\sqrt2$
        apart, and half of that is $\sqrt2$.
      \end{answer}}{\lines{6}}

    \qpart[3]{Now solve the dual problem for the same two points.
      Show that the constraint forces $\alpha_1=\alpha_2=\alpha$, write
      $\mathcal{L}$ as a function of $\alpha$ and maximise it.  Recover
      $\bm{w}$ and check that it agrees with part~(a).  Find $b$ from
      the fact that $\bm{x}_1$ lies on the margin,
      $y_1\,(\bm{w}^\tr\bm{x}_1 - b) = 1$.
      \begin{answer}
        $\alpha_1 y_1 + \alpha_2 y_2 = \alpha_1 - \alpha_2 = 0$, so
        $\alpha_1=\alpha_2=\alpha$.  The inner products are
        $\bm{x}_1^\tr\bm{x}_1 = \bm{x}_2^\tr\bm{x}_2 = 2$ and
        $\bm{x}_1^\tr\bm{x}_2 = -2$, so the double sum is
        \begin{align*}
          \alpha^2\left(2 + 2 + 2\times(+1)(-1)(-2)\right) = 8\,\alpha^2,
          \qquad
          \mathcal{L}(\alpha) = 2\alpha - 4\alpha^2 .
        \end{align*}
        $\mathcal{L}'(\alpha) = 2 - 8\alpha = 0$ gives $\alpha = \tfrac14$.
        Then
        $\bm{w} = \tfrac14(1,1)^\tr - \tfrac14(-1,-1)^\tr = (\tfrac12,\tfrac12)^\tr$,
        as in part~(a).  Finally $\bm{w}^\tr\bm{x}_1 - b = 1 - b = 1$, so
        $b=0$.
      \end{answer}}{\lines{9}}

    \qpart[2]{Add a third point, $\bm{x}_3 = (3,0)^\tr$ with
      $y_3 = +1$.  Without solving the problem again, explain why the
      solution does not change, and why $\alpha_3 = 0$.
      \begin{answer}
        With the old solution, $y_3\,(\bm{w}^\tr\bm{x}_3 - b) = 1.5 > 1$,
        so the new constraint is already satisfied.  Adding a constraint
        that the old optimum satisfies cannot give a better optimum, so
        the solution is unchanged.  Complementary slackness,
        $\alpha_3\left(y_3(\bm{w}^\tr\bm{x}_3 - b) - 1\right) = 0$ with
        the bracket equal to $0.5$, forces $\alpha_3=0$.  Only the
        support vectors (points on the margin) have $\alpha_k>0$.
      \end{answer}}{\lines{6}}

    \qpart[1]{Classify the new point $\bm{x} = (-2,1)^\tr$ using the dual
      form $\mathrm{sgn}\bigl(\sum_k \alpha_k\,y_k\,\bm{x}_k^\tr\bm{x} - b\bigr)$.
      \begin{answer}
        The $\alpha_3$ term vanishes, so the argument is
        $\tfrac14\,(1,1)\,(-2,1)^\tr - \tfrac14\,(-1,-1)\,(-2,1)^\tr
        = -\tfrac14 - \tfrac14 = -\tfrac12$: class $-1$.
      \end{answer}}{\lines{3}}

    \qpart[2]{Return to the two points of part~(a), but use a soft margin
      with $C=0.1$, so that the dual constraint becomes
      $0\leq\alpha_k\leq C$.  Find $\alpha$, $\bm{w}$ and the margin.
      Taking $b=0$, find the slack $s_k = 1 - y_k\,(\bm{w}^\tr\bm{x}_k - b)$
      of each point.
      \begin{answer}
        $\mathcal{L}(\alpha) = 2\alpha - 4\alpha^2$ increases up to
        $\alpha=\tfrac14$, which is now not allowed, so $\alpha = C = 0.1$.
        Then $\bm{w} = 0.1\,(1,1)^\tr + 0.1\,(1,1)^\tr = (0.2,0.2)^\tr$ and
        the margin is $1/\len{\bm{w}} = 1/(0.2\sqrt2) \approx 3.54$, wider
        than before.  Each point has $y_k\,\bm{w}^\tr\bm{x}_k = 0.4$, so
        $s_k = 0.6$: both points sit inside the margin, though they are
        still classified correctly.  A small $C$ makes slack cheap, so
        the SVM buys a wider margin with it.

        (For the curious: here $b$ is not fixed.  Any $|b|\leq0.6$ gives
        $s_1+s_2 = 1.2$ and the same objective.  Symmetry picks $b=0$.)
      \end{answer}}{\lines{8}}
  \end{qparts}
\end{question}
\clearpage

% ***************************** question 4 **************************


\begin{question}{10}
  A function $K(\bm{x},\bm{y})$ is a valid (positive semi-definite)
  kernel if it can be written as
  $K(\bm{x},\bm{y}) = \bm{\phi}(\bm{x})^\tr\bm{\phi}(\bm{y})$ for some
  real feature vector $\bm{\phi}(\bm{x})$.  Equivalently, for any points
  $\bm{x}_1,\ldots,\bm{x}_n$ the Gram matrix with elements
  $K_{ij} = K(\bm{x}_i,\bm{x}_j)$ is positive semi-definite.
  \begin{qparts}
    \qpart[2]{For $\bm{x},\bm{y}\in\mathbb{R}^2$, find a feature vector
      $\bm{\phi}(\bm{x})$ with six components such that
      $(1+\bm{x}^\tr\bm{y})^2 = \bm{\phi}(\bm{x})^\tr\bm{\phi}(\bm{y})$.
      \begin{answer}
        Expanding,
        \begin{align*}
          (1 + x_1y_1 + x_2y_2)^2 = 1 + 2x_1y_1 + 2x_2y_2 + x_1^2y_1^2
          + x_2^2y_2^2 + 2x_1x_2y_1y_2 ,
        \end{align*}
        so $\bm{\phi}(\bm{x}) = (1,\; \sqrt2\,x_1,\; \sqrt2\,x_2,\; x_1^2,\;
        x_2^2,\; \sqrt2\,x_1x_2)^\tr$.  The features are all the
        monomials of degree up to two.
      \end{answer}}{\lines{6}}

    \qpart[2]{Write down the Gram matrix of the linear kernel
      $K(\bm{x},\bm{y}) = \bm{x}^\tr\bm{y}$ for the three points
      $\bm{x}_1 = (1,0)^\tr$, $\bm{x}_2 = (0,1)^\tr$ and
      $\bm{x}_3 = (1,1)^\tr$, and find its eigenvalues.  Why is one of
      them zero?
      \begin{answer}
        \begin{align*}
          \mat{K} = \begin{pmatrix} 1 & 0 & 1\\ 0 & 1 & 1\\ 1 & 1 & 2
          \end{pmatrix}, \qquad
          \det(\mat{K}-\lambda\mat{I}) = -\lambda(\lambda-1)(\lambda-3),
        \end{align*}
        so the eigenvalues are $0$, $1$ and $3$, all non-negative.  The
        zero comes from $\bm{x}_3 = \bm{x}_1 + \bm{x}_2$: then
        $\mat{K}\,(1,1,-1)^\tr = \bm{0}$.  Three points in a
        2-dimensional feature space must be linearly dependent.
      \end{answer}}{\lines{7}}

    \qpart[2]{Show that every valid kernel has $K(\bm{x},\bm{x})\geq0$.
      Hence show that $K(\bm{x},\bm{y}) = \bm{x}^\tr\bm{y} - 1$ is not a
      valid kernel.
      \begin{answer}
        $K(\bm{x},\bm{x}) = \bm{\phi}(\bm{x})^\tr\bm{\phi}(\bm{x}) =
        \len{\bm{\phi}(\bm{x})}^2\geq0$.  But
        $\bm{0}^\tr\bm{0} - 1 = -1 < 0$.
      \end{answer}}{\lines{4}}

    \qpart[2]{Show that $K(\bm{x},\bm{y}) = \len{\bm{x}-\bm{y}}^2$ is not a
      valid kernel, even though $K(\bm{x},\bm{x})\geq0$.  (Hint: look at
      the Gram matrix of two different points.)
      \begin{answer}
        For $\bm{x}\neq\bm{y}$ let $d = \len{\bm{x}-\bm{y}}^2 > 0$.  The
        Gram matrix is
        $\begin{pmatrix} 0 & d\\ d & 0\end{pmatrix}$, with eigenvalues
        $\pm d$.  The eigenvalue $-d$ is negative, so the matrix is not
        positive semi-definite.  (Equivalently, $\bm{v} = (1,-1)^\tr$
        gives $\bm{v}^\tr\mat{K}\bm{v} = -2d < 0$.)
      \end{answer}}{\lines{5}}

    \qpart[2]{Using the rules from the lecture (sums, positive
      multiples, products and exponentials of valid kernels are valid),
      show that
      $K(\bm{x},\bm{y}) = 3\,(\bm{x}^\tr\bm{y})^2 + \exp(\bm{x}^\tr\bm{y})$
      is a valid kernel.
      \begin{answer}
        $\bm{x}^\tr\bm{y}$ is valid, with $\bm{\phi}(\bm{x}) = \bm{x}$.
        Its product with itself, $(\bm{x}^\tr\bm{y})^2$, is valid, and so
        is $3$ times it.  The exponential of a valid kernel,
        $\exp(\bm{x}^\tr\bm{y})$, is valid, and the sum of the two valid
        kernels is valid.
      \end{answer}}{\lines{5}}
  \end{qparts}
\end{question}


\end{document}




